La fédération internationale d'escrime établit des normes de fabrication sur les lames des armes des tireurs afin de garantir au maximum leur sécurité.
	
	
	Pour tester la conformité de l'acier employé, la lame est pliée puis redressée toutes les secondes jusqu'à la rupture.
	
	La lame est conforme si la rupture intervient après plus de cinq heures de test.
	
	\medskip
	
	Un équipementier d'escrime se fournit auprès de trois fabricants de lames.
	
	Son stock est composé de 60~\% de lames du fournisseur A, de 12~\% de lames du fournisseur B, le reste venant du fournisseur C.
	
	\medskip
	
	Une étude qualité a montré que :
	\begin{itemize}[label={\textbullet~}]
		\item 90~\% des lames du fournisseur A étaient conformes;
		\item 95~\% des lames du fournisseur B étaient conformes;
		\item 85~\% des lames du fournisseur C étaient conformes.
	\end{itemize}
	
	\subsection*{Partie A}
	Un contrôle est déclenché sur une des lames vendues par l'équipementier.
	
	\medskip
	
	On considère les évènements suivants :
	\begin{itemize}[label={\textbullet~}]
		\item $A$ : \og La lame testée provient du fournisseur A \fg{};
		\item $B$ : \og La lame testée provient du fournisseur B \fg{};
		\item $C$ : \og La lame testée provient du fournisseur C \fg{};
		\item $T$: \og La lame testée est conforme \fg{}.
	\end{itemize}
	
	On note $\overline{T}$ l'évènement contraire de $T$.
	
	\begin{enumerate}
		\item Recopier et compléter l'arbre de probabilité ci-dessous représentant la situation.
		
		\begin{center}
			
			\begin{tikzpicture}[xscale=1,yscale=1]
				
				\tikzstyle{noeud}=[fill=white,circle,inner sep=2pt]
				\tikzstyle{feuille}=[fill=white,circle,inner sep=2pt]
				\tikzstyle{etiquette}=[midway,fill=white, inner xsep=3pt, inner ysep=1.5pt]
				
				\def\DistanceInterNiveaux{3}
				\def\DistanceInterFeuilles{1}
				
				\def\NiveauA{(0)*\DistanceInterNiveaux}
				\def\NiveauB{(1)*\DistanceInterNiveaux}
				\def\NiveauC{(2)*\DistanceInterNiveaux}
				\def\InterFeuilles{(-1)*\DistanceInterFeuilles}
				
				\node[noeud] (R) at ({\NiveauA},{(2.5)*\InterFeuilles}) {$$};
				\node[noeud] (Ra) at ({\NiveauB},{(0.5)*\InterFeuilles}) {$A$};
				\node[feuille] (Raa) at ({\NiveauC},{(0)*\InterFeuilles}) {$T$};
				\node[feuille] (Rab) at ({\NiveauC},{(1)*\InterFeuilles}) {$\overline{T}$};
				\node[noeud] (Rb) at ({\NiveauB},{(2.5)*\InterFeuilles}) {$B$};
				\node[feuille] (Rba) at ({\NiveauC},{(2)*\InterFeuilles}) {$T$};
				\node[feuille] (Rbb) at ({\NiveauC},{(3)*\InterFeuilles}) {$\overline{T}$};
				\node[noeud] (Rc) at ({\NiveauB},{(4.5)*\InterFeuilles}) {$C$};
				\node[feuille] (Rca) at ({\NiveauC},{(4)*\InterFeuilles}) {$T$};
				\node[feuille] (Rcb) at ({\NiveauC},{(5)*\InterFeuilles}) {$\overline{T}$};
				
				\draw (R.east)--(Ra.west) node[etiquette] {$\dots$};
				\draw (Ra.east)--(Raa.west) node[etiquette] {$\dots$};
				\draw (Ra.east)--(Rab.west) node[etiquette] {$\dots$};
				\draw (R.east)--(Rb.west) node[etiquette] {$\dots$};
				\draw (Rb.east)--(Rba.west) node[etiquette] {$\dots$};
				\draw (Rb.east)--(Rbb.west) node[etiquette] {$\dots$};
				\draw (R.east)--(Rc.west) node[etiquette] {$\dots$};
				\draw (Rc.east)--(Rca.west) node[etiquette] {$\dots$};
				\draw (Rc.east)--(Rcb.west) node[etiquette] {$\dots$};
			\end{tikzpicture}
		\end{center}
		
		\item Déterminer $P\big(A \cap \overline{T}\,\big)$ et interpréter ce résultat dans le contexte de l'exercice.
		
		\item Démontrer que la probabilité que la lame testée soit conforme est de 0,892.
		
		\item Sachant que la lame testée n'est pas conforme, déterminer la probabilité qu'elle provienne du fournisseur B. \emph{On donnera la valeur arrondie au millième.}
	\end{enumerate}
	
	\subsection*{Partie B}
	D'après la \textbf{partie A}, la probabilité qu'une lame de l'équipementier ne soit pas conforme est égale à 0,108.
	
	Lors d'une compétition d'escrime, l'équipementier apporte un échantillon de 75 lames provenant de son stock.
	
	On considère qu'il les a choisies au hasard et de manière indépendante.
	
	De plus son stock de lames est suffisamment grand pour que ce choix puisse être assimilé à un tirage avec remise.
	
	On désigne par $X$ la variable aléatoire égale au nombre de lames non conformes dans cet échantillon.
	
	\begin{enumerate}
		\item On admet que la variable aléatoire $X$ suit une loi binomiale.
		
		Préciser ses paramètres.
		
		\item Calculer la probabilité que 6 lames exactement soient non conformes dans cet échantillon. \emph{On arrondira le résultat au millième.}
		
		\item L'équipementier affirme que la probabilité qu'il y ait strictement plus de 8 lames non conformes dans cet échantillon est inférieure à 50~\%.
		
		A-t-il raison ?
	\end{enumerate}
	
	\subsection*{Partie C}
	Soit $n$ un entier strictement positif désignant le nombre de compétitions durant lesquelles l'équipementier est présent.
	
	Il apporte à chaque fois un échantillon de 75 lames.
	
	Les variables aléatoires $X_{1}$, \ldots, $X_{n}$\,, que l'on considère indépendantes, donnent pour chaque échantillon le nombre de lames non conformes.
	
	\smallskip
	
	On pose \quad $M_{n}=\dfrac{X_{1}+X_{2}+\cdots+X_{n}}{n}$.
	
	\begin{enumerate}
		\item Déterminer l'espérance $\mathrm{E}\left(M_{n}\right)$ et la variance $\mathrm{V}\left(M_{n}\right)$ de la variable aléatoire $M_{n}$\,.
		
		\item Justifier l'inégalité suivante, pour tout entier naturel $n$ non nul :
		\[P\left(\left|\strut M_{n}-8,1\right| \geqslant  2\right) \leqslant  \dfrac{\np{1,8063}}{n}.\]
		
		\item Déterminer une valeur de l'entier $n$ à partir de laquelle \quad $P\left(\left|\strut M_{n}-8,1\right|<2\right) \geqslant  0,95$.
		
		Interpréter ce résultat dans le contexte de l'exercice.
	\end{enumerate}
